% -*- mode: latex; coding: utf-8 -*-
%-----------------------------------------------------------------------------
% pxUML -- User manual and API reference
% Part of P3CTeX. Author: Scvria.
%
% Build from directory tex/ (TEXINPUTS so package is found):
%   set TEXINPUTS=.;./latex;./code;
%   pdflatex -interaction=nonstopmode doc/pxUML.tex
%   makeindex -s gind.ist -o pxUML.ind pxUML.idx
%   pdflatex -interaction=nonstopmode doc/pxUML.tex
%
% Or from tex/doc/: set TEXINPUTS=../latex:../code; then
%   pdflatex -interaction=nonstopmode pxUML.tex
%   makeindex -s gind.ist -o pxUML.ind pxUML.idx
%   pdflatex -interaction=nonstopmode pxUML.tex
%-----------------------------------------------------------------------------
\documentclass[11pt,a4paper]{article}
\usepackage{pxUML}
\usepackage[T1]{fontenc}
\usepackage[utf8]{inputenc}
\usepackage{lmodern}
\usepackage{amsmath}
\usepackage{doc}
\usepackage[colorlinks,urlcolor=blue,linkcolor=blue,hyperindex]{hyperref}

\EnableCrossrefs

\title{The \texttt{pxUML} package\\
  \large UML class diagrams (pgf-umlcd) for P3CTeX / UOC}
\author{Scvria\\P3CTeX / UOC PEC Repository}
\date{2025/02 -- Version 0.1}
\begin{document}
\maketitle
\begin{abstract}
  \texttt{pxUML} provides a LaTeX2e API for UML class and object diagrams on top of
  \texttt{pxpgf-umlcd} (and TikZ). Class definitions are stored in an internal
  registry (via \texttt{pxPRP}); drawing is separate from definition, so the
  same logic can be reused and layout adjusted without changing the model. It is
  tailored for the UOC use case (e.g. PEC diagrams).
\end{abstract}
\tableofcontents

\section{Introduction}
\label{sec:intro}

\texttt{pxUML} keeps a \emph{class registry}: each class has a name, attributes,
operations, inheritance, object type (e.g.\ \texttt{class}, \texttt{interface},
\texttt{enumclass}), visibility, and optional implementation. You define classes
with \texttt{\textbackslash PXdefClass} and optionally extend them with
\texttt{\textbackslash addAttribute}, \texttt{\textbackslash addOperation}, and
\texttt{\textbackslash addInheritance}. Drawing commands (\texttt{\textbackslash drawclass},
\texttt{\textbackslash drawsimpleclass}, \texttt{\textbackslash drawInstance}) read from
this registry, so you can change layout without redefining the class.

\section{Loading the package}
\label{sec:loading}

\begin{verbatim}
\usepackage{pxUML}
\end{verbatim}

The package requires \texttt{expl3}, \texttt{xparse}, \texttt{l3keys2e},
\texttt{color}, \texttt{pxpgf-umlcd}, \texttt{pxPRP}, and \texttt{pgf-umlsd}.

\section{User command reference}
\label{sec:api}

\subsection{Class definition}

\DescribeMacro{\PXdefClass}
Defines or overwrites a class in the registry. The starred form sets visibility
to ``private'' (e.g. for diagram styling).
\begin{verbatim}
\PXdefClass[*][<objtype>]{<ClassName>}[<inheritance>]{<attributes>}[<operations>][<implementation>]
\end{verbatim}
\begin{itemize}
  \item \texttt{<objtype>}: \texttt{class} (default), \texttt{interface},
    \texttt{enumclass}, etc., as in the underlying pgf-umlcd environment.
  \item \texttt{<attributes>}: comma-separated list of attribute strings
    (e.g.\ \texttt{- name : String}).
  \item \texttt{<operations>}: optional; comma-separated list of operation
    strings.
  \item \texttt{<implementation>}: optional; implementation class for interfaces.
\end{itemize}
Example: \texttt{\textbackslash PXdefClass\{MyClass\}\{- id : int\}\{[+ get() : int]\}}.

\subsection{Extending a class}

\DescribeMacro{\addAttribute}
Appends an attribute to the list of attributes of an existing class.
\begin{verbatim}
\addAttribute{<ClassName>}{<attribute>}
\end{verbatim}

\DescribeMacro{\addOperation}
Appends an operation to the list of operations of an existing class.
\begin{verbatim}
\addOperation{<ClassName>}{<operation>}
\end{verbatim}

\DescribeMacro{\addInheritance}
Adds an inheritance (parent) to the class. The second argument is the
\emph{parent class name} only (as expected by the underlying
\texttt{\textbackslash inherit} of pgf-umlcd).
\begin{verbatim}
\addInheritance{<ClassName>}{<ParentClassName>}
\end{verbatim}
Example: \texttt{\textbackslash addInheritance\{DerivedClass\}\{BaseClass\}}.

\subsection{Drawing}

\DescribeMacro{\drawclass}
Draws a full class box (attributes and operations) from the registry at the
given TikZ coordinates. Optional first argument \texttt{[...]} is passed as
TikZ options (e.g.\ \texttt{text width=14em}). Starred form uses the package
``variable'' colour scheme for that class.
\begin{verbatim}
\drawclass[*][<options>]{<ClassName>}[<x,y>]
\end{verbatim}

\DescribeMacro{\drawsimpleclass}
Same as \texttt{\textbackslash drawclass} but draws only the simple class form
(attributes and inheritance, no operations list). Same optional behaviour for
star and options.
\begin{verbatim}
\drawsimpleclass[*][<options>]{<ClassName>}[<x,y>]
\end{verbatim}

\DescribeMacro{\drawInstance}
Draws an object instance (lifeline-style box) of a class, with explicit
attribute and optional operation list.
\begin{verbatim}
\drawInstance[*][<text width em>]{<InstanceName>}[<x,y>]{<ClassName>}{<attr list>}[<op list>]
\end{verbatim}
Default text width is 16\,em. Starred form uses the variable colour scheme.

\DescribeMacro{\PXuml}
Puts the given TikZ/pgf-umlcd content inside a resized box and applies the
package colour scheme. Starred form uses the variable colour scheme for the
whole diagram. Optional width defaults to \texttt{.8\textbackslash linewidth}.
\begin{verbatim}
\PXuml[*][<tikz options>]{<content>}[<width>]
\end{verbatim}
Example: \texttt{\textbackslash PXuml\{\textbackslash drawclass\{MyClass\}\}[1\textbackslash linewidth]}.

\subsection{Inspection}

\DescribeMacro{\PXshowclass}
Loads the stored data for the class into internal token lists (for debugging or
further macros). If the class does not exist in the registry, all keys are
effectively empty; no error is raised.
\begin{verbatim}
\PXshowclass{<ClassName>}
\end{verbatim}

\section{Minimal example}
\label{sec:example}

\begin{verbatim}
\usepackage{pxUML}
\begin{document}
\PXdefClass{Person}{- name : String}[+ getName() : String]
\addAttribute{Person}{- age : Integer}

\PXdefClass{Base}{- id : int}
\PXdefClass{Derived}{- extra : String}
\addInheritance{Derived}{Base}

\PXuml{
  \drawclass[text width=12em]{Person}[0,0]
  \drawsimpleclass{Base}[6,0]
  \drawsimpleclass{Derived}[6,-3]
}[0.9\linewidth]
\end{document}
\end{verbatim}

\section{Change history}
\label{sec:changes}

\begin{description}
  \item[v0.1] Initial release: class registry, \texttt{\textbackslash PXdefClass},
    \texttt{\textbackslash drawclass}, \texttt{\textbackslash drawsimpleclass},
    \texttt{\textbackslash drawInstance}, \texttt{\textbackslash PXuml},
    \texttt{\textbackslash PXshowclass}, \texttt{\textbackslash addAttribute},
    \texttt{\textbackslash addOperation}, \texttt{\textbackslash addInheritance}.
    Documentation and tests aligned with pxPRP workflow.
\end{description}

\PrintIndex
\end{document}
